% --- Learning outcomes ---
\begin{frame}{Learning Outcomes}
  By the end of this lecture you should be able to:
  \begin{itemize}\itemsep0.4em
    \item[\textcolor{red}{$\blacktriangleright$}] \textbf{Compute} the memory footprint of a training run --- parameters, gradients, optimizer states, activations --- and predict when and why it will run out of memory.
    \item[\textcolor{red}{$\blacktriangleright$}] \textbf{Explain and choose between} data, tensor, pipeline, and sequence parallelism, and say what each one costs in communication.
    \item[\textcolor{red}{$\blacktriangleright$}] \textbf{Describe} ZeRO stages 1--3 and FSDP sharding, and configure them for a fine-tuning run.
    \item[\textcolor{red}{$\blacktriangleright$}] \textbf{Combine} strategies into a 3D-parallel plan for a large pre-training job on a given cluster.
    \item[\textcolor{red}{$\blacktriangleright$}] \textbf{Explain} how distributed \emph{inference} differs from distributed \emph{training}: replication versus sharding, routing, and failure handling.
  \end{itemize}
  \vspace{0.5em}
  \begin{itemize}\footnotesize
    \item[\textcolor{red}{$\blacktriangleright$}] \textbf{Assumed background} (all covered earlier): the transformer block, scaling laws, and single-node inference optimisation --- KV caching, continuous batching, quantization.
  \end{itemize}
\end{frame}
